\documentclass[tikz]{standalone}
\usepackage{geometricfigures}

\usepackage{amsmath,amssymb}
\usepackage{xcolor}
\usetikzlibrary{calc,arrows.meta,positioning}

\begin{document}
\begin{tikzpicture}[
                remember picture, % Enable referencing nodes from this tikzpicture later
                transform shape
            ]

                % Position it below the identity, roughly aligned
                \coordinate (zero_point) at (0, 0);
                \node[above of=zero_point, yshift=-2em] {$\mathbb{O}$};
                \draw[thick, muted, fill=muted!20!bg, opacity=0.8] ($(zero_point) + (-1.0, 0.0)$) -- ($(zero_point) + (2.5, -0.0)$) -- ($(zero_point) + (1.5, -2.5)$) -- ($(zero_point) + (-2.0, -2.5)$) -- cycle;
                \node[muted] at ($(zero_point) + (0.5, -2.2)$) {$\mathfrak{g} \simeq T_\mathbb{I} G$};

                % Example element in Lie Algebra
                \coordinate (arrow_tip) at ($(zero_point) + (0.5, -1.5)$);
                \draw[->, thick, fg] (zero_point) -- (arrow_tip) node[midway, xshift=1.5em] {$X \in \mathfrak{g}$};
\end{tikzpicture}

\end{document}
